\documentclass[11pt,a4paper]{article}

\usepackage[utf8]{inputenc}
\usepackage[T1]{fontenc}
\usepackage[french]{babel}
\usepackage{amsmath,amssymb}
\usepackage[top=1.2cm,bottom=1.5cm,left=2cm,right=2cm]{geometry}
\usepackage{graphicx}
\usepackage{booktabs}
\usepackage{hyperref}
\hypersetup{colorlinks=true,linkcolor=black,citecolor=blue,urlcolor=blue}

\title{Question C.3 --- Invariance du budget de preuve}
\author{}
\date{}

\begin{document}
\maketitle

\noindent Toutes les mesures numériques de cette section proviennent de la campagne
instrumentée de 2\,000 exécutions (20 amplitudes $\times$ 100 graines, $H=2\,000$,
socle $\hat p_0$ estimé sur 3\,000 pas pré-rupture) documentée dans
\texttt{04\_experimentations/ulysse/}. La hiérarchie déterministe
$\tau_{\mathrm{ARF}} \le \tau_{\mathrm{swap}}(q)$ (Q~C.1) et l'analyse dimensionnelle
avec son certificat de détection (Q~C.2) sont supposées acquises.

\section{Expression du budget et exposant résiduel}

Le modèle stylisé de l'article suppose qu'après la rupture l'erreur excède son socle
d'une quantité $\Delta e$ pendant un transitoire de longueur $W = \tau_{\mathrm{ARF}}$,
puis retombe au socle. L'aire d'erreur excédentaire offerte au détecteur vaut alors
$\mathbb{E}[A] = \Delta e \cdot \mathbb{E}[\tau_{\mathrm{ARF}}]$, à condition que
l'horizon couvre entièrement le transitoire, ce que $H = 2\,000$ garantit sur toute la
grille --- l'amplitude la plus faible donne $3 \times 18{,}5/0{,}028 \approx 1\,982$. En y
injectant l'ajustement empirique de l'article
$\mathbb{E}[\tau_{\mathrm{ARF}}] \approx 18{,}5\,\Delta e^{-0,98}$~:
\begin{equation}
  \boxed{\mathbb{E}[A] \approx 18{,}5\,\Delta e^{\,1-0,98} = 18{,}5\,\Delta e^{\,0,02}}
  \label{eq:budget}
\end{equation}

\paragraph{Exposant résiduel et variation sur la grille.} L'exposant résiduel de
$\Delta e$ vaut $0{,}02$. Sur $\Delta e \in [0{,}028\,;\,0{,}498]$, le rapport des
extrêmes est $(0{,}498/0{,}028)^{0,02} = 1{,}059$~: le budget prédit varie de $5{,}9\%$,
de $17{,}2$ à $18{,}2$, quand l'amplitude du choc varie d'un facteur $17{,}7$.

\paragraph{Conclusion formelle.} L'amplitude et la durée se compensent presque
exactement : la forêt se répare d'autant plus vite que le choc est fort, et le produit
$\Delta e \times W$ reste quasi constant. Un choc macroéconomique violent ne laisse donc
pas plus de preuve au détecteur qu'une dérive lente. Le budget se comporte comme une
constante du couple forêt-détecteur, et non comme une caractéristique du drift observé.

\section{Le budget réellement utilisable}

Le détecteur ne consomme pas l'aire brute. Sa tolérance $\delta_P = 0{,}01$ est
retranchée à chaque pas, si bien que le budget effectivement accumulable vaut
$(\Delta e - \delta_P)\,W$ et non $\Delta e\,W$. La part confisquée est donc
$\delta_P/\Delta e$, indépendante de la durée du transitoire.

\begin{center}
\begin{tabular}{lccc}
\toprule
$\Delta e$ & budget brut & part prélevée & budget utilisable \\
\midrule
$0{,}028$ (bas de grille) & $17{,}23$ & $\mathbf{35{,}5\%}$ & $11{,}11$ \\
$0{,}498$ (haut de grille) & $18{,}24$ & $2{,}0\%$ & $17{,}88$ \\
\bottomrule
\end{tabular}
\end{center}

\paragraph{Où la tolérance mord.} C'est à l'extrémité \emph{basse} de la grille que
$\delta_P$ prélève une part non négligeable : plus d'un tiers du budget à
$\Delta e = 0{,}028$, contre $2\%$ à l'autre extrémité. La tolérance étant un seuil
absolu, elle pèse d'autant plus lourd que le signal est faible. Elle aggrave donc
l'invariance au lieu de la corriger : le budget utilisable n'est plus constant, il
décroît vers les faibles amplitudes --- là où le détecteur en aurait le plus besoin.

\section{Confrontation au certificat en fenêtre entière}

Le certificat établi en~Q~C.2 exige, sur la fenêtre complète, une aire brute d'au moins
$\lambda + H\delta_P$. Avec $H = 2\,000$ et $\delta_P = 0{,}01$ le péage vaut $20$ quel
que soit le seuil, d'où $28$, $45$ et $70$ pour $\lambda \in \{8, 25, 50\}$, à comparer
au budget d'environ $18$ établi en~\eqref{eq:budget}. \textbf{Aucun des trois seuils
n'est atteignable}, le plus bas exigeant déjà une fois et demie le budget disponible.
La preuve requise n'existe pas dans le signal, quel que soit le réglage retenu.

\paragraph{Ce que devient le certificat quand l'horizon s'allonge.} Le seuil requis
$\lambda + H\delta_P$ croît \emph{linéairement} en $H$, tandis que le budget reste
\emph{constant} d'après~\eqref{eq:budget}. L'écart se creuse donc indéfiniment : une fois
l'erreur résorbée, chaque pas d'observation supplémentaire ajoute $\delta_P$ au péage
sans rien apporter au numérateur, si bien qu'observer plus longtemps dégrade le
certificat au lieu de l'améliorer. D'où l'intérêt de l'ancrer sur une fenêtre courte
$w = 3 \times 18{,}5/\Delta e$, taillée sur le transitoire de l'amplitude considérée.

\section{Vérification sur les mesures}

\begin{figure}[htbp]
  \centering
  \includegraphics[width=0.8\textwidth]{ulysse/resultats/figures/Fig_QC_budget_full.png}
  \caption{Aire d'erreur excédentaire mesurée en fonction de $\Delta e$, campagne
  complète, comparée à la constante prédite par~\eqref{eq:budget}. En bleu l'aire sur la
  fenêtre courte $A(w)$, en rouge l'aire sur l'horizon entier $A(H)$.}
  \label{fig:budget}
\end{figure}

\paragraph{L'invariance tient dans son sens, pas dans son chiffre.} Sur
$\Delta e \in [0{,}085\,;\,0{,}498]$, l'amplitude du choc varie d'un facteur $5{,}8$, et
sans adaptation de la forêt l'aire varierait d'autant. Elle varie en réalité d'un facteur
$\mathbf{1{,}78}$~: l'effet compensatoire est bien présent. Mais le modèle prédit un
facteur $1{,}04$, et il sous-estime le budget au milieu de la grille, avec $30{,}8$
mesuré contre $18{,}0$ prédit à $\Delta e = 0{,}243$. L'invariance est donc confirmée
dans son principe et approximative dans ses valeurs.

\paragraph{L'aire en fenêtre entière devient négative.} $A(H)$ culmine à $38{,}3$ puis
décroît, passe sous zéro à partir de $\Delta e = 0{,}482$ et descend jusqu'à $-20{,}8$.
L'explication tient à la géométrie du problème plutôt qu'à la mesure. À grande amplitude
la frontière post-rupture est $x_0 + x_1 > 4$ avec $x_0 + x_1 \sim \mathcal{N}(0,2)$~:
la classe 1 devient quasi absente et la tâche post-drift se trouve être plus facile que
la tâche initiale. L'erreur mesurée tombe de $0{,}024$ avant la rupture à $0{,}007$ en
fin d'horizon, et ce déficit soustrait une trentaine d'unités à l'aire brute sur les
$1\,900$ pas suivant la résorption. Le modèle stylisé est donc en défaut à forte
amplitude, où l'erreur passe sous son socle initial au lieu d'y retomber. C'est une
seconde raison, indépendante de celle donnée en~Q~C.2, de ne certifier que sur la
fenêtre courte.

\paragraph{D'où vient l'écart : le socle, la censure, ou l'ajustement ?} Le sujet invite
à regarder d'abord du côté de $\hat p_0$. Son erreur d'estimation se propage de façon
parfaitement corrélée sur toute l'aire, avec un biais $-H(\hat p_0 - p_0)$~; estimé sur
3\,000 pas, $\hat p_0$ a un écart-type inter-graines de $0{,}00298$, soit $\pm 6{,}0$ sur
$A(H)$. C'est non négligeable devant un budget de $18$ --- mais insuffisant pour rendre
compte de l'écart de $14{,}5$ constaté à $\Delta e = 0{,}028$, d'autant que ce biais est
de signe aléatoire alors que l'écart observé est systématique.

Le test qui départage vraiment consiste à comparer $\tau_{\mathrm{ARF}}$ mesuré à sa
propre prédiction, puisque \eqref{eq:budget} superpose deux hypothèses distinctes~: le
modèle $\mathbb{E}[A] = \Delta e\,\mathbb{E}[\tau_{\mathrm{ARF}}]$ d'une part, et
l'ajustement en loi de puissance d'autre part.

\begin{center}
\begin{tabular}{lrrrrr}
\toprule
$\Delta e$ & $\tau_{\mathrm{ARF}}$ mesuré & prédit & ratio & $\Delta e \times \tau_{\mathrm{ARF}}$ & $A(w)$ mesuré \\
\midrule
$0{,}028$ & $118$ & $611$ & $\mathbf{0{,}19}$ & $3{,}33$ & $2{,}71$ \\
$0{,}141$ & $263$ & $126$ & $\mathbf{2{,}08}$ & $37{,}07$ & $27{,}09$ \\
$0{,}287$ & $64$ & $63$ & $1{,}02$ & $18{,}37$ & $29{,}22$ \\
$0{,}436$ & $34$ & $42$ & $0{,}81$ & $14{,}82$ & $23{,}16$ \\
$0{,}498$ & $28$ & $37$ & $0{,}76$ & $13{,}93$ & $17{,}31$ \\
\bottomrule
\end{tabular}
\end{center}

À la plus faible amplitude, l'ajustement prédit le premier remplacement au pas $611$
alors qu'il survient au pas $118$, cinq fois plus tôt. Le budget recalculé à partir du
$\tau_{\mathrm{ARF}}$ mesuré vaut $3{,}33$, ce qui reproduit l'aire observée de $2{,}71$.
Le modèle $\mathbb{E}[A] = \Delta e\,\mathbb{E}[\tau_{\mathrm{ARF}}]$ tient donc --- c'est
la loi de puissance qui est en cause. L'écart n'est pas localisé : le rapport
mesuré sur prédit vaut $0{,}19$ en bas de grille, dépasse $2$ vers
$\Delta e = 0{,}141$, ne colle qu'autour de $\Delta e = 0{,}287$, puis retombe à $0{,}76$
en haut.

\begin{center}
\fbox{\parbox{0.93\textwidth}{\textbf{Conséquence pour le rapport.} L'invariance du
budget repose sur deux hypothèses successives. La première, le modèle
$\mathbb{E}[A] = \Delta e \cdot \mathbb{E}[\tau_{\mathrm{ARF}}]$, est confirmée par les
mesures. La seconde, l'injection de $18{,}5\,\Delta e^{-0,98}$, ne l'est qu'au voisinage
du milieu de la grille. La conclusion « un choc violent ne laisse pas plus de preuve
qu'une dérive lente » reste valide, mais elle ne peut plus être justifiée par l'exposant
$0{,}02$~: il faut l'appuyer sur la mesure directe de $\tau_{\mathrm{ARF}}$.}}
\end{center}

\paragraph{Le socle théorique.} La cible pré-rupture $y = \mathbf{1}\{x_0 + x_1 > 0\}$
est une fonction \emph{déterministe} des entrées, sans bruit d'étiquetage : l'erreur de
Bayes vaut exactement zéro. Le socle mesuré, $0{,}0231$, est donc intégralement de
l'erreur d'approximation --- les arbres de Hoeffding découpant par coupes perpendiculaires
aux axes une frontière oblique. Le détecteur surveille ainsi les écarts à un niveau de
référence qui dépend du classifieur lui-même : une forêt mieux ajustée abaisserait le
socle et modifierait le budget de preuve sans qu'aucune propriété du drift n'ait changé.
%
% Note : le passage de 1000 a 3000 pas pour estimer p_0 divise le bruit par 1,68
% (attendu sqrt(3)=1,73). Les conclusions agregees ne bougent pas -- mediane de tau_det
% inchangee a 21 pas pour lambda=8 -- car les deux estimateurs sont non biaises, seule
% leur variance differe. Le gain porte sur la mesure par execution, donc sur les
% correlations de Q D.

\section{Ce que cette section établit}

Le budget de preuve offert au détecteur est quasi indépendant de l'amplitude du drift :
l'aire varie d'un facteur $1{,}8$ quand l'amplitude varie d'un facteur $5{,}8$, là où
elle suivrait cette dernière en l'absence d'adaptation. La tolérance $\delta_P$ ampute
ce budget aux faibles amplitudes --- et allonger l'horizon d'observation creuse encore
le déficit.

Deux réserves accompagnent ce résultat. L'ajustement en loi de puissance sur lequel
l'article fonde l'invariance n'est valide qu'au voisinage du milieu de la grille, si bien
que l'argument doit reposer sur la mesure directe de $\tau_{\mathrm{ARF}}$ plutôt que sur
l'exposant $0{,}02$. Et le modèle stylisé lui-même est en défaut à forte amplitude, où
l'erreur passe sous son socle initial.
%
% TODO(a synchroniser) : cette section fournit les trois nombres (28, 45, 70) que Q C.2
% confronte a son certificat, et la constante d'invariance que Q D mobilise pour
% expliquer l'effondrement du tau_b global par agregation. Verifier la coherence des
% renvois une fois les trois redactions reunies.

\end{document}
